\section{Elevator Rides}
\label{sec:cses-elevator-rides}

\problemstatement{There are $n$ people with given weights, and an elevator
with maximum weight $x$. Find the minimum number of elevator rides needed
to carry everyone (people may be grouped arbitrarily, but each ride's
total weight must not exceed $x$).}

\begin{patternbox}
With $n\le20$, ``partition people into fewest groups under a capacity''
screams \textbf{bitmask DP}: the state is the \emph{subset} of people
already transported. The clever part is what to store -- not just the ride
count, but how full the last ride is.
\end{patternbox}

\subsection*{A two-part state per subset}

For each subset \textit{mask} of people, store a pair
$\textit{dp}[\textit{mask}]=(\textit{rides},\textit{lastWeight})$: the
minimum number of rides to carry exactly that subset, and among solutions
with that few rides, the smallest weight currently loaded in the last
(still-open) ride. Compare pairs lexicographically -- fewer rides first,
then a lighter last ride (which leaves the most room). Transition by adding
one new person $p$ to $\textit{mask}$:
\begin{itemize}
  \item If $p$ fits in the current last ride
        ($\textit{lastWeight}+w_p\le x$): same ride count, last weight
        grows by $w_p$.
  \item Otherwise start a new ride: rides $+1$, last weight $=w_p$.
\end{itemize}

\begin{ideabox}[Store Rides \emph{and} the Last Ride's Load]
Minimising rides alone is not enough, because two subsets with equal ride
counts differ in how much room remains -- and that room affects future
additions. Pairing the ride count with the least-loaded last ride makes
the greedy-within-DP choice optimal.
\end{ideabox}

\subsection*{Algorithm}

\begin{enumerate}
  \item $\textit{dp}[0]=(1,0)$ -- one empty ride, zero load.
  \item For each mask and each person $p\notin\textit{mask}$: compute
        the candidate pair for $\textit{mask}\cup\{p\}$ using the rule
        above; keep the lexicographically smaller.
  \item Answer is the ride count of $\textit{dp}[\text{full mask}]$.
\end{enumerate}

\subsection*{Dry run}

\dryrun{Weights $[2,3,3]$, capacity $x=5$.}

\begin{itemize}
  \item $\textit{dp}[\{0\}]=(1,2)$, $\textit{dp}[\{1\}]=(1,3)$.
  \item Add person $2$ ($w=3$) to $\{0\}$ (last $2$): $2+3=5\le5$, same
        ride: $(1,5)$. To $\{1\}$ (last $3$): $3+3=6>5$, new ride:
        $(2,3)$.
  \item $\textit{dp}[\{0,2\}]=(1,5)$, then adding person $1$
        ($3$): $5+3>5$, new ride $(2,3)$. Full-set answer: $2$ rides
        (e.g.\ $\{2,3\}$ and $\{3\}$).
\end{itemize}

\subsection*{C++ implementation}

\begin{lstlisting}
#include <bits/stdc++.h>
using namespace std;

int main() {
    int n, x;
    cin >> n >> x;
    vector<int> w(n);
    for (int& v : w) cin >> v;

    int full = 1 << n;
    // dp[mask] = {rides, weight in last ride}
    vector<pair<int,int>> dp(full, {n + 1, 0});
    dp[0] = {1, 0};

    for (int mask = 1; mask < full; mask++) {
        for (int p = 0; p < n; p++) {
            if (!(mask & (1 << p))) continue;
            auto prev = dp[mask ^ (1 << p)];       // state without p
            if (prev.second + w[p] <= x)           // fits in last ride
                prev.second += w[p];
            else                                   // needs a new ride
                prev = {prev.first + 1, w[p]};
            dp[mask] = min(dp[mask], prev);        // lexicographic
        }
    }
    cout << dp[full - 1].first << "\n";
    return 0;
}
\end{lstlisting}

\begin{complexitybox}
\textbf{Time Complexity.} $O(2^n\cdot n)$ -- each of $2^n$ subsets tries
adding each of $n$ people. \textbf{Space Complexity.} $O(2^n)$ for the DP
table.
\end{complexitybox}

\begin{takeawaybox}
Bitmask DP fits partition problems with $n\le20$. The lesson beyond the
mask is the \emph{compound state}: pairing ``rides so far'' with ``room
left in the current ride'' captures exactly the extra information a greedy
packing needs -- a recurring trick when a count alone loses too much.
\end{takeawaybox}
